\section{Materiais e Métodos}\label{sec:metodos}

\begin{figure*}[t]
  \centering
  \resizebox{\textwidth}{!}{% Diagramas de blocos de confiabilidade (RBD) dos três motores. Blocos empilhados: par redundante.
\begin{tikzpicture}[
  font=\scriptsize,
  blk/.style={draw, rounded corners=1pt, align=center, minimum height=6.6mm, inner sep=1.5pt, text width=#1},
  blk/.default=12mm,
  rot/.style={font=\footnotesize\bfseries, anchor=east},
  >=stealth, line width=0.4pt]
\def\g{2.2mm}
% ---------------- Pistão ----------------
\node[rot] (lp) at (0,0) {Pistão};
\node[blk=11mm, right=\g of lp] (p1) {Admissão};
\coordinate[right=\g of p1] (p2i);
\node[blk=11mm, right=\g of p2i, yshift=4.6mm] (p2a) {Bomba mecânica};
\node[blk=11mm, right=\g of p2i, yshift=-4.6mm] (p2b) {Bomba elétrica};
\coordinate (p2o) at ($(p2a.east |- p2i)+(\g,0)$);
\node[blk=13mm, right=\g of p2o] (p3) {Carburador / injeção};
\coordinate[right=\g of p3] (p4i);
\node[blk=11mm, right=\g of p4i, yshift=4.6mm] (p4a) {Magneto 1};
\node[blk=11mm, right=\g of p4i, yshift=-4.6mm] (p4b) {Magneto 2};
\coordinate (p4o) at ($(p4a.east |- p4i)+(\g,0)$);
\node[blk=19mm, right=\g of p4o] (p5) {Cilindros, válvulas, pistões, virabrequim};
\node[blk=13mm, right=\g of p5] (p6) {Lubrificação};
\node[blk=15.5mm, right=\g of p6] (p7) {Arrefecimento};
\node[blk=13.5mm, right=\g of p7] (p8) {Hélice e governador};
\draw (p1) -- (p2i) |- (p2a); \draw (p2i) |- (p2b);
\draw (p2a) -| (p2o); \draw (p2b) -| (p2o); \draw (p2o) -- (p3);
\draw (p3) -- (p4i) |- (p4a); \draw (p4i) |- (p4b);
\draw (p4a) -| (p4o); \draw (p4b) -| (p4o); \draw (p4o) -- (p5);
\draw (p5) -- (p6) -- (p7) -- (p8);
% ---------------- Turboélice ----------------
\node[rot] (lt) at (0,-1.75) {Turboélice};
\node[blk=10mm, right=\g of lt] (t1) {Entrada de ar};
\node[blk=13mm, right=\g of t1] (t2) {Compressor};
\node[blk=13mm, right=\g of t2] (t3) {Câmara de combustão};
\node[blk=17mm, right=\g of t3] (t4) {Turbinas (compressor e livre)};
\node[blk=11mm, right=\g of t4] (t5) {Caixa de redução};
\node[blk=13.5mm, right=\g of t5] (t6) {Hélice e governador};
\node[blk=13mm, right=\g of t6] (t7) {Combustível (FCU)};
\node[blk=13mm, right=\g of t7] (t8) {Lubrificação};
\draw (t1) -- (t2) -- (t3) -- (t4) -- (t5) -- (t6) -- (t7) -- (t8);
% ---------------- Turbofan ----------------
\node[rot] (lf) at (0,-3.35) {Turbofan};
\node[blk=7mm, right=\g of lf] (f1) {Fan};
\node[blk=15.5mm, right=\g of f1] (f2) {Compressores LP e HP};
\node[blk=13mm, right=\g of f2] (f3) {Câmara de combustão};
\node[blk=12mm, right=\g of f3] (f4) {Turbinas HP e LP};
\coordinate[right=\g of f4] (f5i);
\node[blk=11mm, right=\g of f5i, yshift=4.6mm] (f5a) {FADEC A};
\node[blk=11mm, right=\g of f5i, yshift=-4.6mm] (f5b) {FADEC B};
\coordinate (f5o) at ($(f5a.east |- f5i)+(\g,0)$);
\node[blk=13mm, right=\g of f5o] (f6) {Combustível};
\node[blk=13mm, right=\g of f6] (f7) {Lubrificação};
\node[blk=12mm, right=\g of f7] (f8) {Caixa de acessórios};
\draw (f1) -- (f2) -- (f3) -- (f4) -- (f5i) |- (f5a); \draw (f5i) |- (f5b);
\draw (f5a) -| (f5o); \draw (f5b) -| (f5o); \draw (f5o) -- (f6) -- (f7) -- (f8);
\end{tikzpicture}
}
  \caption{Diagramas de blocos de confiabilidade dos três motores. Blocos empilhados formam um par redundante, que só falha se as duas unidades falharem.}
  \label{fig:rbd}
\end{figure*}

\subsection{Motores estudados}\label{sec:motores}

Cada tecnologia é representada por um modelo de uso difundido: o Lycoming O-540 (pistão), o Pratt \& Whitney Canada PT6A-34 (turboélice) e o CFM56-7B (turbofan).
O O-540 tem seis cilindros horizontais opostos, refrigerados a ar, e segue o ciclo de quatro tempos.
A cada volta, dezenas de peças aceleram, param e invertem o sentido, o que as expõe a fadiga e desgaste~\cite{faa_amt}.
A ignição é duplicada: cada cilindro tem duas velas, alimentadas por dois magnetos que geram a própria energia e não dependem da bateria~\cite{faa_phak}.
Na alimentação, a bomba acionada pelo motor é complementada por uma bomba elétrica auxiliar.
Os demais subsistemas não têm redundância.

O PT6A-34 é um turboélice de turbina livre.
O gerador de gás retira dos gases apenas a energia que gira o compressor, e a turbina livre, sem ligação mecânica com ele, aciona a hélice por uma caixa de redução~\cite{faa_amt}.
A unidade de controle de combustível (FCU) dosa o combustível, e um governador ajusta o passo da hélice.
A ignição só atua na partida e fica fora do modelo; o turboélice não tem blocos redundantes.
O CFM56-7B, que equipa o Boeing 737 de nova geração, tem dois eixos concêntricos: o de baixa pressão liga o fan e o compressor de baixa à turbina de baixa, e o de alta liga o compressor de alta à turbina de alta.
O controle é feito por um FADEC (\emph{full authority digital engine control}) de dois canais, cada um capaz de controlar o motor sozinho~\cite{faa_amt}.

\subsection{Diagramas de blocos}\label{sec:rbd}

A Fig.~\ref{fig:rbd} mostra os RBD.
O evento de falha é a perda de potência que obriga a desligar o motor ou o impede de entregar a potência necessária ao voo.
Todo subsistema cuja falha, sozinha, provoca esse evento entra em série.
Em relação à proposta, as duas turbinas do turboélice formam agora um único bloco, porque os códigos do NTSB não as distinguem; com blocos em série e taxa constante, isso não altera $\lambda$, MTBF ou $R(t)$.
Ficam fora os sistemas da aeronave, como os tanques, e componentes como escapamento, partida e indicação.

Supõe-se taxa de falha constante e independência entre blocos, de modo que $R_i(t)=e^{-\lambda_i t}$ e, em série, $R_s=\prod_i R_i$.
Nos pares redundantes, uma fração $\beta$ da taxa $\lambda_u$ de cada unidade corresponde a causas comuns, como combustível contaminado nas duas bombas ou erro de manutenção repetido nos dois magnetos, que derrubam as duas unidades ao mesmo tempo~\cite{ebeling1997}:
\begin{equation}
  R_p(t) = \left[1-\left(1-e^{-(1-\beta)\lambda_u t}\right)^2\right] e^{-\beta\lambda_u t}.
  \label{eq:beta}
\end{equation}
A confiabilidade do motor é o produto de $R_s$ e dos $R_p$, e $\beta$ varia de 0,01 a 0,3, com valor nominal de 0,1.
Com pares, $R(t)$ deixa de ser exponencial.
Por isso, o MTBF sem reparo é calculado como a integral de $R(t)$, e o MTBF com o par restaurado a cada voo, como o inverso da taxa equivalente $-\ln R(1\,\mathrm{h})/1\,\mathrm{h}$; o segundo caso corresponde ao cheque de magnetos antes da decolagem.
Para dois motores na mesma aeronave, aplicamos~\eqref{eq:beta} com o motor no lugar da unidade e um $\beta$ entre motores, que representa causas como combustível contaminado, manutenção feita pela mesma equipe e gelo.
Nesse caso, o evento é a perda total de propulsão.

\subsection{Taxas totais}\label{sec:taxas}

\begin{table}[t]
  \caption{Taxas de falha publicadas (falhas \pcmh)}
  \label{tab:taxas}
  \centering
  \begin{tabular}{@{}lll@{}}
    \toprule
    Motor ou exigência & Taxa & Fonte \\
    \midrule
    Pistão Lycoming, aviões leves (2009--2014) & 12,7 & \cite{atsb2016} \\
    Pistão Continental, aviões leves (2009--2014) & 12,1 & \cite{atsb2016} \\
    Pistão, estimativa da FAA & 10 a 100 & \cite{romeyn2007} \\
    Turboélice PT6 (2016) & $\approx$0,15 & \cite{stobie2019} \\
    Turbina, exigência SET-IMC & $<1$ & \cite{easa2017} \\
    Turbofan CFM56-7B (fabricante) & 0,20 & \cite{cfm2012} \\
    Turbofan, limite ETOPS de 180 min & $\le 2$ & \cite{cfr25k} \\
    \bottomrule
  \end{tabular}
\end{table}

A taxa total de cada motor vem da literatura (Tabela~\ref{tab:taxas}), com mínimo, nominal e máximo de 10, 12,4 e 100 \pcmh{} para o pistão, 0,15, 0,15 e 1 para o turboélice e 0,10, 0,20 e 2 para o turbofan.
O nominal do pistão é a média de Lycoming e Continental~\cite{atsb2016}, e o do turbofan vem de um comunicado da CFM~\cite{cfm2012}, sem estatística independente da FAA ou da EASA.
Os limites regulatórios entram só como máximos, e o mínimo do turbofan é assumido.

A taxa publicada $\lambda$ é preservada em um voo de 1~h: os blocos em série recebem $\lambda_s=\lambda-\sum_p[-\ln R_p(1\,\mathrm{h})]$, isto é, o que sobra depois da contribuição dos pares.
Uma consequência é que $\beta$ e $\lambda_u$ apenas redistribuem a taxa entre os blocos, sem alterar a probabilidade de falha em 1~h.

\subsection{Base do NTSB}\label{sec:ntsb}

A única base de ocorrências é a do NTSB~\cite{ntsb}, com as tabelas de eventos, aeronaves, motores, sequência de eventos e achados (\emph{findings}) de 2008 a 2026.
O universo são 18\,735 aviões-eventos com motor a pistão, turboélice ou turbofan, sem os 3258 de construção amadora, porque as taxas da literatura são de motores certificados.
Há falha de motor quando a sequência de eventos registra perda total ou parcial de potência, mau funcionamento do grupo motopropulsor, falha não contida ou corte do motor.
Se a mesma sequência tem um evento de combustível (pane seca, falta de alimentação, contaminação ou combustível errado), a perda é operacional e é contada à parte.
São 2562 falhas no pistão, 176 no turboélice e 161 no turbofan.

Os achados codificam o componente por seção e subseção, que seguem os capítulos ATA.
Usamos os achados de aeronave dos grupos hélice e motopropulsor dos aviões com falha de motor e atribuímos cada código ao bloco pelo prefixo mais longo de uma tabela de correspondência (por exemplo, 8530 cilindros, núcleo; 7414 magneto; 7250 turbinas).
O código de controle eletrônico de combustível (7321) vai para o FADEC no turbofan e para a FCU no turboélice.
Ficam fora os erros de uso pelo piloto (77 achados), os códigos genéricos que não identificam o bloco, como \emph{engine (reciprocating)} (535), e os componentes sem bloco no RBD (146); cada avião conta uma vez por bloco.
Ao todo, 1109 registros foram atribuídos a blocos: 994 do pistão (880 nos blocos em série), 69 do turboélice e 46 do turbofan (45 em série).

\subsection{Frações por subsistema}\label{sec:fracoes}

A taxa $\lambda_s$ é distribuída entre os blocos em série por
\begin{equation}
  \lambda_i = \lambda_s \, \frac{n_i}{N_s},
  \label{eq:fracao}
\end{equation}
em que $n_i$ é o número de registros atribuídos ao bloco $i$ e $N_s$ é o total nos blocos em série.
O intervalo de 95\% de cada fração é o de Jeffreys, isto é, os quantis de $\mathrm{Beta}(n_i+\tfrac12,\,N_s-n_i+\tfrac12)$~\cite{brown2001}; ele reflete apenas a incerteza amostral, e não o viés da fonte.
A taxa de cada unidade redundante foi estimada na mesma escala.
Cada registro de bloco em série vale $\lambda/N_s$, e o registro de um par se refere a uma das duas unidades, de modo que $\lambda_u = n_p\lambda/(2N_s)$, com $n_p$ registros do par.
Um par sem registro recebe a média da posteriori de Jeffreys no lugar de $n_p/(n_p+N_s)$, porque nenhum registro não significa taxa nula.

\subsection{Tempo de motor e consequência da falha}\label{sec:tempo}

O NTSB informa as horas do motor desde novo (TSN) e desde a última revisão geral (TSO), que usamos só nos monomotores, porque nos multimotores não se sabe qual motor falhou.
Sem horas voadas da frota, testamos a taxa constante por um desenho caso-controle.
Os casos são os monomotores com falha de motor, e os controles são os monomotores acidentados por outra causa, sem evento de motor, sem evento de combustível e sem achado do grupo motopropulsor; supõe-se que a idade do motor dos controles represente a da frota.
Se a taxa de falha é Weibull, $h(t)\propto t^{k-1}$, a chance de um avião com motor de idade $t$ ser caso cresce com $t^{k-1}$, e a regressão logística
\begin{equation}
  \operatorname{logit} P(\text{caso}\mid t) = a + b\ln t
  \label{eq:caso}
\end{equation}
dá o parâmetro de forma $k=1+b$, com intervalo de Wald; $k=1$ é a taxa constante do RBD.
A mesma base dá a proporção de falhas de motor com vítima fatal, por tecnologia e número de motores, com intervalo de Wilson~\cite{wilson1927}.

\subsection{Rede bayesiana e atualização das frações}\label{sec:bayes}

O RBD foi convertido em uma rede bayesiana equivalente~\cite{bobbio2001}, com um nó-raiz por bloco em série, três por par (unidades A e B e causa comum) e nós determinísticos até o nó do motor.
Com no máximo 12 raízes, a distribuição conjunta é enumerada e a inferência é exata.
A rede responde a duas perguntas: dada a falha do motor, qual bloco é a causa mais provável; e quanto sobe a probabilidade de falha quando uma unidade de um par já está inoperante.
As frações também foram atualizadas por um modelo Dirichlet--multinomial, com as frações iniciais da proposta como \emph{prior} (20 pseudocontagens).
A proposta previa ainda atualizar as taxas absolutas por um modelo Gama--Poisson, mas o NTSB não traz horas voadas por frota; sem a exposição, essa etapa não foi executada.

\subsection{Sensibilidade e implementação}\label{sec:sensib}

Os índices foram calculados para um voo de 1~h e para um ano de 1000~h, e a sensibilidade leva cada parâmetro ($\lambda$, $\beta$ e $\lambda_u$) do mínimo ao máximo de sua faixa, com os demais no valor nominal.
Todo o processamento foi feito em Python, em um \emph{pipeline} reprodutível que gera as tabelas e figuras deste artigo a partir dos dados brutos do NTSB.

\subsection{Declaração de uso de inteligência artificial}\label{sec:ia}

\emph{Ferramenta:} Claude Code, assistente de programação da Anthropic, com o modelo Claude Opus.
\emph{Finalidade:} escrever e depurar o código Python do \emph{pipeline} (tratamento das bases, contagens, modelo RBD, rede bayesiana, estatística e figuras) e apoiar a redação e a formatação em \LaTeX{} deste artigo.
\emph{Estratégia:} a equipe definiu o escopo, as hipóteses de modelagem e os critérios de classificação, e montou a correspondência entre os códigos de componente do NTSB e os blocos do RBD; a ferramenta produziu código e rascunhos de texto a partir dessas decisões e dos resultados; todo número do texto vem das tabelas geradas pelo \emph{pipeline}, e as referências foram conferidas nas fontes originais.
\emph{Responsabilidade:} os autores revisaram o código, os resultados e o texto, e respondem integralmente pelo conteúdo do trabalho.
